% answers_short.tex —— 1997 年·普通高等学校招生考试·全国卷文科 · 简短答案速查
% 由 tools/make_short_answers.py 自 answers.tex 的 \gksolhead/\gkanswer 提取，
% 并与 1997/1997_数学真题_全国文_答案要点.md 逐题比对后生成，勿手改（改源 MD 后重跑）。

\begin{gknotice}
\textbf{使用说明}：本节只给最终答案便于快速核对，完整推导见解析版。答案为编写者独立推导，\textbf{非官方标准答案}。
\end{gknotice}

\gkpart{选择题}
\begin{center}\begin{tabular}{cccccccccccccccc}
\toprule
题号 & 1 & 2 & 3 & 4 & 5 & 6 & 7 & 8 & 9 & 10 & 11 & 12 & 13 & 14 & 15 \\
\midrule
答案 & B & B & A & C & B & C & D & C & A & B & A & D & C & C & B \\
\bottomrule
\end{tabular}\end{center}

\gkpart{填空题}
\begin{itemize}[leftmargin=2.6em]
  \item[\textbf{16.}] 4
  \item[\textbf{17.}] $(4,2)$
  \item[\textbf{18.}] $2-\sqrt3$
  \item[\textbf{19.}] ①、④
\end{itemize}

\gkpart{解答题}
\begin{itemize}[leftmargin=2.6em]
  \item[\textbf{20.}] 模 $\sqrt2$；辐角主值 $5\pi/6$
  \item[\textbf{21.}] $a_n=4-\frac{12}{5}(n-1)$ 或 $a_n=-3+\frac{16}{5}(n-1)$
  \item[\textbf{22.}] (1) $y=S(bv+a/v)$，$0<v\le c$；(2) $v=\min(c,\sqrt{a/b})$
  \item[\textbf{23.}] (1) $AD\perp D_1F$；(2) $90^\circ$；(3) 未决；(4) 1
  \item[\textbf{24.}] (1) $O,C,D$ 共线；(2) $A=(\sqrt3,\frac16\log_2 3)$
  \item[\textbf{25.}] $(x-1)^2+(y-1)^2=2$ 或 $(x+1)^2+(y+1)^2=2$
\end{itemize}
